\documentclass[11pt,twocolumn]{article}

\usepackage[utf8]{inputenc}
\usepackage[T1]{fontenc}
\usepackage{amsmath,amssymb}
\usepackage{natbib}
\usepackage{hyperref}
\usepackage{geometry}
\usepackage{booktabs}
\usepackage{graphicx}
\usepackage{enumitem}
\usepackage{microtype}

\geometry{margin=1in}

\hypersetup{
  colorlinks=true,
  linkcolor=blue!60!black,
  citecolor=blue!60!black,
  urlcolor=blue!60!black
}

\title{\textbf{Emergence: A Programming Paradigm for\\Constraint-Shaped Agency}}

\author{
  Jeremy McEntire\\
  \texttt{jeremy@wander.com}
}

\date{February 2026}

\begin{document}

\maketitle

\begin{abstract}
The programmer does not write the program.  The programmer writes the
conditions under which the program emerges.  We propose \emph{Emergence}
as a programming paradigm in which the fundamental unit of computation
is an autonomous agent that interprets messages within a constraint
architecture---schemas, thresholds, scoring functions, and topological
rules---and the system's behavior arises from the interaction of agents
operating within those constraints.  This paradigm recovers Alan Kay's
original vision of object-oriented programming, which required a
computational substrate with contextual understanding that did not exist
until large language model agents.  We ground the paradigm in convergent
results from cybernetics, modular design theory, organizational theory,
and complex systems.  We present a working stigmergic mesh as an
existence proof and address the paradigm's real limitations---latency,
cost, non-determinism, and the genuinely new discipline of designing
constraint spaces---with equal honesty.
\end{abstract}

%% ============================================================
\section{Introduction}
\label{sec:intro}

Every programming paradigm answers one question: \emph{where does the
intelligence live?}  In procedural programming, it lives in the sequence
of instructions.  In functional programming, it lives in the composition
of transformations.  In object-oriented programming---as practiced, not
as conceived---it lives in the class hierarchy.

In each case, the programmer encodes the intelligence directly.  The
system does what it is told.  This works when the problem is fully
specifiable in advance.  It fails when the problem space exceeds the
programmer's capacity to enumerate its states---which is to say, it
fails for every interesting problem in every sufficiently complex
domain.

The arrival of large language model (LLM) agents creates a genuinely
new possibility.  For the first time, the computational unit itself
possesses contextual understanding.  An LLM agent assessing a code
review comment does not pattern-match against a list of phrases; it
understands what code review is \emph{for}, recognizes when a comment is
substantive versus ceremonial, and distinguishes a genuine technical
concern from a territorial objection.  This capability is not
incremental improvement.  It is a qualitative change in what a
computational unit \emph{can be.}

We call the resulting paradigm \textbf{Emergence}.  The taxonomy is
clean:

\begin{quote}
\emph{Procedural programming: you write the steps.
Object-oriented programming: you write the entities.
Functional programming: you write the transformations.
Emergence: you write the constraints, and the behavior writes
itself.}
\end{quote}

The program is not the behavior.  The program is the physics, the
constraints, the container, the field equations.  The behavior is the
emergent property of the system the programmer designed---not the system
itself.  This shift implies not merely a new tool but a new discipline:
the skill of designing constraint spaces that reliably produce useful
behavior is neither programming as currently practiced nor management as
currently understood, and its development is as much a part of the
paradigm's contribution as its technical architecture
(Section~\ref{sec:between}).

%% ============================================================
\section{The Kay Recovery}
\label{sec:kay}

Alan Kay has spent decades clarifying what he meant---and what he did
not mean---by ``object-oriented programming.''  The clarifications have
been largely ignored.

In a 2003 email, Kay wrote: ``I'm sorry that I long ago coined the term
`objects' for this topic because it gets many people to focus on the
lesser idea.  The big idea is messaging'' \citep{kay2003meaning}.  In
his Smalltalk work and subsequent reflections, Kay consistently
described three properties:

\begin{enumerate}[nosep]
  \item Each object is an autonomous entity---``a recursion of the whole
    computer''---with its own state, behavior, and interpretation of
    incoming messages.
  \item Communication happens exclusively through messages, and the
    \emph{receiver} decides what a message means.
  \item The system's behavior emerges from the interaction of these
    autonomous message-interpreting entities, not from a central
    controller.
\end{enumerate}

What the industry built was different.  Method dispatch replaced message
interpretation: a message became a function call, resolved at compile
time or through a vtable, with the caller knowing exactly what would
happen.  Class hierarchies replaced autonomy: objects became instances
of types, their behavior determined by inheritance chains.  Design
patterns replaced emergence: the Gang of Four catalog
\citep{gamma1994design} codified strategies for managing complexity that
the paradigm was supposed to dissolve.

The misimplementation was not stupidity.  It was a technological
constraint.  In 1972---or 1995, or 2015---there was no computational
substrate that could \emph{interpret} a message.  A function can parse
one.  A pattern matcher can classify one.  A rule engine can apply
conditions to one.  But none of these \emph{understands} one.  The
receiver cannot decide what a message means if the receiver has no
capacity for meaning.

Kay's vision required a computational unit with genuine contextual
understanding.  It required exactly what an LLM agent provides.

%% ============================================================
\section{The Agent as Computational Primitive}
\label{sec:agent}

An LLM agent is not a better function.  It is a different kind of
computational entity.  The distinction matters because it is the
distinction Kay was trying to draw.

A function \emph{computes}: given inputs, it produces outputs according
to a fixed algorithm.  An agent \emph{assesses}: given a situation, it
produces a judgment shaped by its understanding and context.  A function
that detects scope creep by counting changed lines is doing something
fundamentally different from an agent that reads a pull request and
recognizes that what was filed as a ``bug fix'' is actually a feature
rewrite touching authentication, database schema, and three API
contracts simultaneously.  The function applies a rule.  The agent
understands a situation.

This distinction enables three properties previously unavailable:

\paragraph{Genuine message interpretation.}  When an agent receives a
signal---a pull request, a customer ticket, a deployment log---it does
not merely parse the text.  It interprets meaning in context.  A
large change touching authentication logic receives different attention
than an identically sized change adding stylesheets, not because of a
rule about filenames, but because the agent understands what
authentication code is and why changes to it carry risk.

\paragraph{Judgment under ambiguity.}  Traditional computation requires
the programmer to anticipate every case.  An agent exercises judgment in
situations the programmer never considered, because its training
encompasses the breadth of human knowledge about how to reason in
context.  This is the same capability that allows a new employee to
handle situations not covered by the employee handbook.

\paragraph{Compositional emergence.}  When multiple agents, each
interpreting messages within their own context, interact through a
shared environment, the system-level behavior is not the sum of the
agents' behaviors.  Patterns, correlations, and structures arise that no
individual agent was designed to detect.  This is the property that Kay
described and that class-hierarchy OOP could never deliver.

%% ============================================================
\section{The Constraint Architecture}
\label{sec:constraints}

If the agents provide the flexibility, what provides the reliability?
This is the central engineering question of the paradigm, and its
answer is the paradigm's central insight: \textbf{the program is the
boundary, not the behavior.}

In Emergence, the programmer designs the constraint architecture within
which agents operate: \emph{rigid at the interfaces, flexible in the
interiors.}

\subsection{Rigid Boundaries}

The boundaries are non-negotiable structural elements that ensure
coherence:

\begin{description}[nosep,leftmargin=1em]
  \item[Schemas.] The data structures agents consume and produce are
    strictly typed.  A signal has a source, timestamp, content, and
    metadata.  An assessment has a confidence score, a set of claims, and
    a provenance chain.  The agent decides what to \emph{put} in these
    fields; the schema decides what fields \emph{exist}.
  \item[Thresholds.] Numerical parameters that gate behavior without
    prescribing it.  A vigilance parameter determines how similar a
    signal must be to existing knowledge before an agent accepts it.  An
    energy budget determines how many signals an agent can process before
    it must yield.  The agent exercises judgment \emph{within} these
    bounds; the bounds themselves are fixed.
  \item[Scoring functions.] Mathematical functions that evaluate outputs
    without understanding them.  A familiarity score combines lexical
    overlap, semantic similarity, structural alignment, and temporal
    recency into a single number.  The agent does not know it is being
    scored.  The scoring function does not know what the agent is
    thinking.  The decoupling is the point: the meta-system shapes
    selection pressure without micromanaging cognition.
  \item[Topological rules.] Rules governing how agents connect,
    communicate, and compete.  Which agents see which signals.  How
    routing priority is determined.  When agents fork, merge, or decay.
    These rules define the \emph{structure} of interaction without
    constraining the \emph{content}.
\end{description}

\subsection{Flexible Interiors}

Within the boundaries, agents operate with genuine autonomy.  They
interpret signals, form assessments, develop specializations, and
generate insights without procedural instruction.  The programmer does
not specify \emph{how} an agent should assess whether a team's
communication patterns are degrading.  The programmer specifies
\emph{that} communication health is a property to assess, provides a
schema for reporting it, and trusts the agent's contextual
understanding to do the rest.

This is the inversion that defines the paradigm.  In traditional
programming, the interior is rigid (algorithms, control flow, explicit
logic) and the boundaries are flexible (APIs change, schemas evolve).
In Emergence, the boundaries are rigid and the interiors are flexible.
The programmer's job shifts from \emph{writing behavior} to
\emph{shaping the space in which behavior occurs.}

\subsection{Recursive Constraint Interaction}

The constraint architecture is not a thermostat with one feedback loop.
It has many degrees of freedom (one per agent per context) and many
feedback loops (scoring, routing, lifecycle, consensus) that
\emph{interact}.  The vigilance threshold affects which signals an agent
accepts, which affects its specialization, which affects the topology,
which affects which signals it receives, which affects whether the
vigilance threshold should change.

This recursive interaction is the mechanism of emergence.  The
constraint architecture creates conditions for self-organization; the
self-organization produces capabilities that no individual constraint
was designed to produce.

\subsection{Well-Formedness}

If the paradigm is to be engineering rather than aspiration, constraint
architectures must have well-formedness conditions.  A constraint
architecture is well-formed when it satisfies three properties:

\begin{enumerate}[nosep]
  \item \textbf{Completeness}: every agent output passes through at
    least one scoring function.  No agent can produce assessments that
    bypass evaluation.  This is the analog of type safety: the system
    guarantees that all behavior is bounded, even if the specific
    behavior is not predicted.
  \item \textbf{Liveness}: the lifecycle rules guarantee that the
    system cannot freeze.  There must exist conditions under which new
    agents are created (fork), idle agents are removed (decay), and
    redundant agents are consolidated (merge).  A system that can only
    grow or only shrink is not well-formed.
  \item \textbf{Dissipation}: the energy budget is finite and
    monotonically decreasing absent external input.  No agent can
    process signals indefinitely.  This prevents monopoly (one agent
    capturing all signals) and guarantees that the system must
    continuously earn its computational expenditure.
\end{enumerate}

These conditions do not specify what the system will do.  They specify
what the system \emph{cannot} do: produce unscored outputs, freeze into
a fixed topology, or consume resources without bound.  The formal
development of these conditions into a full design calculus---constraint
typing, compositional well-formedness, convergence guarantees---is
future work.  But the conditions themselves are already enforced in the
existence proof and are sufficient to distinguish a well-formed
constraint architecture from an ad hoc collection of agents.

%% ============================================================
\section{Theoretical Convergence}
\label{sec:theory}

The Emergence paradigm converges with results from four independent
traditions.  The convergence is not coincidental---each tradition
discovered, through its own methods, the same structural truth about
where rigidity and flexibility must live.

\subsection{Ashby's Law of Requisite Variety}

Ashby's Law \citep{ashby1956introduction} states that a regulator must
have at least as much variety as the system it regulates.  But Ashby's
formulation is more specific than the popular summary suggests.  The law
constrains the \emph{channel capacity} between regulator and
regulated---the bandwidth through which the regulator observes and acts
upon the system.  In traditional programming, the channel capacity is
determined by the programmer's ability to enumerate states.  Every
\texttt{if} branch, every \texttt{case} statement, every validation
rule is an explicit expansion of the channel.  This works until the
state space exceeds enumerability.

Emergence achieves requisite variety through a different channel.  The
agents' contextual understanding \emph{already encompasses} the variety
of the problem space.  The constraint architecture does not need to
enumerate states; it needs to \emph{bound the agents that are already
navigating them.}  The channel capacity comes from the agents, not the
program.  This is why the paradigm applies precisely where traditional
programming fails: in domains where the state space exceeds any
programmer's capacity to enumerate---organizations, ecosystems,
markets, complex social systems.

\subsection{Beer's Viable System Model}

Stafford Beer's Viable System Model (VSM) \citep{beer1972brain}
identifies the 3--4 homeostat as the central tension in any viable
system: System~3 (internal stability) must be dynamically balanced
against System~4 (environmental intelligence).  Beer's specific
contribution is not the observation that stability and adaptability
must coexist---that is platitude.  It is that \emph{System~4 must be a
genuine model of the environment, not a summary of internal reports.}
Organizations fail, Beer argues, because System~4 degrades into a
mirror of System~3: the ``intelligence'' function reports what
operations already believe rather than what the environment actually
contains.

This is precisely the failure mode that Emergence is designed to avoid.
The agents in the flexible interior are not summarizing internal state;
they are interpreting environmental signals with contextual understanding
independent of the system's prior beliefs.  The lifecycle
rules---fork, decay, merge---are the homeostat, and the rigid
boundaries are System~3.  But the specific contribution of Emergence to
Beer's framework is a System~4 that cannot be captured by System~3,
because the agents' understanding comes from outside the system's own
training data.

\subsection{Baldwin and Clark's Design Rules}

Baldwin and Clark \citep{baldwin2000design} demonstrated that modular
systems outperform integral systems when design rules---the interfaces
between modules---are well-specified.  Their framework goes further than
the popular summary ``modularity is good.''  They provide a formal
option-pricing model showing that modularity creates \emph{option
value}: each module can be independently varied, creating a
combinatorial space of possible designs that can be explored in
parallel.  The value of this optionality grows with the number of
modules and the uncertainty of the design space.

In Emergence, agents are modules and constraints are design rules.  Each
agent independently develops specialization within the rigid interface
definitions.  The system explores the space of possible specializations
in parallel---an agent that absorbs deployment signals develops
deployment expertise while another, absorbing customer tickets, develops
support expertise.  Neither was programmed to specialize.  The option
value is not financial but epistemic: the system discovers what
specializations are useful rather than requiring the programmer to
predict them.

\subsection{Kauffman's Edge of Chaos}

Kauffman's work on Boolean networks \citep{kauffman1993origins}
identified the regime between order and chaos where adaptive
computation occurs.  In frozen networks, perturbations die out and the
system cannot adapt.  In chaotic networks, perturbations propagate
everywhere and the system cannot maintain structure.  At the edge of
chaos, perturbations propagate just far enough to enable adaptation
without dissolution.

Kauffman showed that the position in this spectrum is controlled by two
parameters: the connectivity of the network (how many other elements
each element influences) and the bias of the Boolean functions (how
likely each element is to be in a given state).  The Emergence
constraint architecture controls both: topological rules govern
connectivity (which agents see which signals), and thresholds govern
bias (how likely an agent is to accept, reject, or ignore).  The
lifecycle rules---fork, merge, decay---continuously adjust these
parameters, maintaining the system at the edge without requiring the
programmer to calculate the critical regime directly.

\subsection{The Convergence}

Four traditions, developed independently over six decades, arrived at
the same prescription: rigid constraints at the boundaries, flexible
behavior in the interiors, with feedback loops that maintain the
balance.  Ashby specified the information-theoretic requirement.  Beer
described the organizational architecture.  Baldwin and Clark proved the
economic value.  Kauffman identified the computational regime.  Even
Postel's Law---``be conservative in what you send, liberal in what you
accept'' \citep{postel1980dod}---is the same insight applied to protocol
design: rigid outputs, flexible inputs.  Emergence generalizes this
principle from protocols to entire systems: rigid boundaries, flexible
agents.

The convergence across independent formalisms is the strongest evidence
that the paradigm captures something real---not a framework imposed on
recalcitrant material, but a pattern the material itself exhibits.

%% ============================================================
\section{Related Work}
\label{sec:related}

The Emergence paradigm did not arise in a vacuum.  Several research
traditions have explored adjacent territory, and the differences are as
instructive as the similarities.

\paragraph{BDI Agents.}  The Belief-Desire-Intention architecture
\citep{rao1995bdi} gave agents internal mental states---beliefs about
the world, desires about outcomes, intentions about plans.  This was a
genuine advance over reactive agents, but the beliefs, desires, and
intentions were still programmer-specified data structures.  A BDI agent
does not \emph{understand} its beliefs; it \emph{stores} them.  The
architecture provided the right structure but lacked the substrate: a
computational unit that could actually form beliefs from evidence rather
than retrieving them from a database.

\paragraph{Swarm Intelligence.}  Ant colony optimization, particle
swarm optimization, and related approaches
\citep{bonabeau1999swarm} demonstrated that simple agents following
simple rules can produce complex collective behavior.  This is genuine
emergence, and it inspired the stigmergic coordination in our
existence proof.  But swarm agents are \emph{simple}---they follow
fixed behavioral rules without contextual understanding.  Emergence
extends the principle to agents with genuine interpretive capacity,
enabling coordination in domains (language, judgment, assessment) where
simple rules are insufficient.

\paragraph{Autonomic Computing.}  IBM's autonomic computing initiative
\citep{kephart2003vision} proposed self-managing systems organized
around the MAPE-K loop (Monitor, Analyze, Plan, Execute, Knowledge).
The architecture anticipated many elements of Emergence: self-organization,
feedback loops, adaptive behavior.  But the MAPE-K loop assumes that
the Analyze and Plan phases can be implemented as traditional
algorithms.  In domains where analysis requires judgment rather than
computation, the loop stalls---the system can monitor, but it cannot
\emph{understand} what it monitors.

\paragraph{Self-Adaptive Systems.}  The SEAMS community has explored
systems that modify their own behavior at runtime, typically through
models@runtime and feedback control
\citep{cheng2009software}.  These systems are the closest
precursor to Emergence, but they rely on explicit models of the
adaptation space---the programmer must specify what can change and how.
Emergence removes this requirement: the agents' contextual
understanding replaces the explicit adaptation model.

In each case, the limitation was the same: the agents lacked
understanding.  The architectures were right.  The substrate was wrong.
Emergence inherits the architectural insights of these traditions and
pairs them with the first computational substrate capable of delivering
on their promise.

%% ============================================================
\section{The Chronicler Principle}
\label{sec:chronicler}

The paradigm is easiest to understand through analogy.

Consider a system designed to write a novel.  The traditional approach
programs a plot generator: act structure, character arcs, conflict
escalation, prose templates.  The programmer encodes the intelligence
of storytelling into explicit rules.

The Emergence approach builds a \emph{world}: a physics model that
enforces conservation laws, an economy model that enforces scarcity, a
character model that enforces psychological consistency, a geography
model that enforces spatial constraints.  Each model is an agent bound
to a context, interpreting events within its domain.  No model knows
about stories.  No model has a plot.  But when a character model reports
that a character's motivation has shifted, and the economy model reports
that a resource has become scarce, and the geography model reports that
two factions now share a border---a story emerges from the interaction.

A chronicler agent observes the interactions, recognizes narrative
structure in the convergence of model outputs, and produces prose.  The
chronicler does not invent events.  It \emph{witnesses} them.  The
story is not written; it is \emph{discovered.}

This is not metaphor.  It is the literal architecture.  Replace
``physics model'' with ``schema validator'' and ``character model'' with
``domain expert agent'' and ``chronicler'' with ``consensus layer,'' and
you have the architecture of any Emergence system.  The programmer
builds the world.  The agents inhabit it.  The behavior is what happens
when the world's rules interact.

%% ============================================================
\section{Existence Proof: A Stigmergic Mesh}
\label{sec:proof}

Theory without implementation is speculation.  We have built a working
system that instantiates the Emergence paradigm: a stigmergic mesh for
ambient structure discovery in organizations
\citep{mcentire2026asd}.

The system ingests work artifacts---pull requests, issue tickets,
deployment logs, chat messages---as signals.  No one formulates a query.
No one specifies what to look for.  The system processes the ambient
exhaust of organizational work and surfaces structural patterns that no
one asked about.

\subsection{Architecture}

The mesh consists of worker agents, each bound to a context representing
a region of specialization.  Signals arrive and are routed via
breadth-first search through the worker topology.  The first worker
whose familiarity score exceeds a vigilance threshold accepts the
signal; the remaining workers never see it.  Accepted signals update the
worker's context through one of three learning modes: full storage
(indexed), compressed summary, or lossy weight shift.  Workers that
accept too many signals fork.  Workers that accept too few decay.
Redundant workers merge.  The topology self-organizes.

The rigid boundaries are strict: immutable typed signals, assessment
schemas with required confidence scores, a five-component familiarity
function, exponential energy decay with activity-based recharge.  The
flexible interiors are genuine: when an agent assesses a pull request, it
brings the full weight of its contextual understanding to bear on what
the change means, what risks it carries, and what patterns it
participates in.

\subsection{What Emerged}

During live deployment against a real engineering organization
(dozens of repositories, multiple chat channels, several project
tracking teams), the system produced detections that no individual
component was designed to generate:

\begin{itemize}[nosep]
  \item A policy intervention flagged a pull request that modified
    security-sensitive code buried inside a large feature branch---not
    because of a rule about filenames, but because the agent recognized
    that a sensitive surface was being modified incidentally, inside a
    change scoped as something else, and that similar modifications
    were accumulating faster than the budgeted rate.
  \item A churn detector surfaced that a repository had exceeded its
    weekly diff budget, flagging not the volume itself but the
    \emph{pattern}: a single large PR pushing cumulative churn past a
    threshold, suggesting scope creep in what had been filed as a
    contained change.
  \item Workers self-organized into specializations that no one
    configured.  Over successive runs, individual workers developed
    affinities for infrastructure concerns, product-facing behavior,
    and cross-team coordination---not because they were assigned these
    domains, but because the signals they accepted reinforced
    familiarity in those directions.  Workers that received no signals
    decayed and were pruned.  Workers that received too many forked,
    creating new specialists.
  \item Cross-source correlations emerged.  The same organizational
    pattern---a team under sustained pressure---manifested differently
    across version control (shorter review comments, faster merges),
    project tracking (tickets reclassified from feature work to
    defects), and chat (reduced channel activity).  No individual
    source showed the pattern clearly.  The convergence across
    sources---detected by independently specialized agents whose
    outputs were correlated by attention mechanisms---surfaced it.
\end{itemize}

The pattern across these detections is consistent: \emph{temporal
accumulation invisible to any single actor.}  Each individual action---a
pull request, a code review, a ticket reclassification---is visible to
its author and unremarkable in isolation.  The trajectory across
actions, across authors, across sources, is visible to no one.  That is
the mesh's niche: it detects the structural patterns that arise from the
interaction of individually reasonable decisions, which is precisely the
class of problem that traditional query-based systems cannot address
because no one knows to ask.

Equally important is what the system \emph{suppressed.}  The majority
of ingested signals were filtered as retrieval rather than discovery:
factual restatements of what an individual author did, adding no
structural insight.  The filtering machinery---bot detection, duplicate
suppression, retrieval/discovery classification---is not auxiliary to the
emergence; it is what makes the emergence legible.  Without it, the
genuine structural findings would be buried in noise and no one would
read any of them.

The system operates with over one thousand passing tests and several
automated detectors for early indicators of organizational dysfunction.
A full ingestion cycle costs under one dollar and completes in minutes.

%% ============================================================
\section{What Changes}
\label{sec:changes}

\subsection{The Programmer's Role}

The programmer becomes an architect of constraint spaces.  The skills
shift from algorithm design to boundary design: what schemas to enforce,
what thresholds to set, what scoring functions to use, what topological
rules to apply.  The question changes from ``what should the system
do?'' to ``what should the system \emph{not} be able to do?''---and
then trusting that useful behavior emerges within those limits.

This is not deskilling.  Designing constraint architectures that produce
useful emergent behavior is at least as difficult as writing explicit
algorithms.  The difficulty shifts from ``how do I make the system
do~$X$?'' to ``how do I create conditions under which $X$ naturally
arises?''

\subsection{Testing and Verification}

Emergence systems are not fully deterministic.  You do not get the same
novel twice.  You do not get the same mesh findings twice.  You get
\emph{structurally consistent} results---bounded by the same
constraints, shaped by the same scoring functions, governed by the same
lifecycle rules---but the specific trajectories vary because the agents
are flexible.

This is a feature, not a defect---but only if the constraints are
designed well enough that every trajectory is acceptable.  That is the
engineering skill: not writing behavior, but designing the space of
acceptable behaviors so thoroughly that whatever emerges within it is
useful.  This is what makes Emergence engineering rather than hope.

Testing must therefore shift from input-output verification to
\emph{property verification}: does the system maintain its invariants?
Do the constraints hold?  Does the topology self-organize into useful
specializations?  Do cross-agent correlations converge on genuine
patterns rather than noise?  The discipline is not less rigorous; it is
differently rigorous.

\subsection{Failure Modes}

A traditional system fails by producing wrong outputs.  An Emergence
system fails by producing \emph{incoherent} outputs---agents that
contradict each other, emergent patterns that are artifacts rather than
structure, consensus that converges on noise.

Debugging shifts accordingly.  Where traditional debugging traces
execution paths, Emergence debugging examines constraint interactions:
which boundaries are too tight (suppressing useful behavior), which are
too loose (permitting incoherent behavior), which feedback loops are
amplifying noise.  In the stigmergic mesh, this means monitoring energy
landscapes, tracking spectral properties of the agent topology, and
examining whether worker specializations are converging or drifting.

\subsection{The Configuration Trap Dissolved}

A pervasive pattern in traditional software is the \emph{configuration
trap}: flexibility implemented as configuration options accumulates
until the combinatorial space of configurations becomes the dominant
source of complexity.

Emergence dissolves this trap by placing flexibility in the agents
rather than in configuration.  The constraints are few and
well-defined.  The flexibility comes from contextual interpretation,
which does not require configuration---it requires understanding, which
the agents already possess.

%% ============================================================
\section{Limitations}
\label{sec:limitations}

The Emergence paradigm has real limitations that must be acknowledged
honestly.  Some are fundamental; some are contingent on the current
state of technology.

\subsection{Latency}

Every agent assessment is an LLM inference call.  In the stigmergic
mesh, processing a single signal through reflection, evaluation, and
assessment takes seconds, not microseconds.  A backfill of several dozen signals
from a version control system takes several minutes.  This is not a network optimization
problem; it is a fundamental property of the computational substrate.
LLM inference is slow in the same way that human judgment is slow---the
richness of contextual interpretation has a temporal cost.

This means Emergence systems operate on a different timescale than
traditional software.  They are not suitable for real-time control,
latency-sensitive APIs, or tight feedback loops.  They are suitable for
problems where the cost of missing a pattern exceeds the cost of
detecting it minutes or hours later---organizational intelligence,
risk monitoring, strategic assessment.  The paradigm does not replace
traditional computing; it occupies a different niche.

\subsection{Cost}

Agent assessments cost money.  Each LLM call consumes tokens, and
tokens have a price.  The stigmergic mesh spends on the order of one dollar
per session processing dozens of signals across multiple sources.  At
continuous operation, this is a meaningful operational cost.  The
constraint architecture must therefore include explicit economic
boundaries: token budgets, energy decay functions that limit processing,
three-tier learning modes that range from full analysis (expensive) to
lossy impression (cheap).  Cost management is not an operational
afterthought; it is a first-class constraint in the architecture.

\subsection{The Space Between: A New Discipline}
\label{sec:between}

The hardest limitation is human, not technical.  Emergence requires a
skill that does not yet have a name, a training pipeline, or a
professional identity.  It is not programming---the practitioner does
not write algorithms or control flow.  It is not management---the
practitioner does not direct agents through conversation or feedback.
It is something in between: closer to designing a game's rules than to
playing it, closer to landscape architecture than to construction,
closer to creating an ecosystem than to building a machine.

The practitioner must think in terms of selection pressures rather than
instructions, in terms of fitness landscapes rather than control flow,
in terms of what behaviors a constraint \emph{space} permits rather than
what a specific agent will \emph{do.}  The feedback loops are slower---you
cannot step through a debugger; you observe emergent behavior across
many signals and many cycles.  The diagnostic intuition is ecological
rather than mechanical: ``the system isn't specializing because the
vigilance threshold is too low'' is a different kind of reasoning than
``the function returns null because the input is empty.''

This is arguably the paradigm's most important implication.  Every
previous paradigm shift changed what programmers \emph{wrote}---structured
blocks instead of gotos, objects instead of procedures, pure functions
instead of mutable state.  Emergence changes what programmers \emph{are.}
The practitioner becomes something closer to a constitutional
designer---defining the rules under which autonomous entities operate,
the invariants that must hold regardless of what those entities do, and
the feedback mechanisms that keep the system within its viable envelope.
The analogy to governance is not decorative: the questions are the same.
What freedoms do the agents need?  What constraints prevent pathological
outcomes?  How do you design accountability without destroying autonomy?

This discipline will develop.  Training pipelines will emerge.  But in
the near term, the paradigm asks practitioners to operate in a
genuinely unfamiliar cognitive mode---the space between writing code and
herding cats, where the work product is neither an algorithm nor a
directive but a set of conditions under which useful behavior reliably
arises.

\subsection{Non-Determinism}

Emergence systems do not produce identical outputs from identical
inputs.  This requires a precise distinction: the \emph{constraints} are
deterministic; the \emph{trajectories} are not.  The scoring functions
produce the same scores for the same inputs.  The energy budgets decay
at the same rates.  The schemas enforce the same shapes.  What varies is
the path the agents take through the space those constraints define.

This means the debugging methodology is different but not absent: you
debug the constraints, not the behavior.  If the system produces
pathological output, the question is not ``what did agent~$A$ do
wrong?'' but ``what constraint failed to prevent this trajectory?''
The constraints guarantee what \emph{will not} happen.  The agents
discover what \emph{will.}

For many engineering contexts, non-deterministic trajectories are
disqualifying.  The paradigm does not apply where exact reproducibility
is required.  It applies where \emph{structural consistency} is
sufficient: where the question is not ``did the system produce
output~$X$?'' but ``did the system produce an output within the
acceptable set?''  This is a genuine restriction on the paradigm's
scope, not a limitation to be engineered around.

\subsection{State and Memory}

Emergence systems must learn, but LLM agents are stateless between
calls.  All continuity must be externalized---into context summaries,
weight vectors, learned term lists, communication graphs, finding
registries.  This externalization is lossy.  The stigmergic mesh
preserves semantic state (what agents have learned) but loses
topological state (how many agents exist, their specialization profiles,
their neighbor relationships) on restart.  The system reconverges, but
the reconvergence path may differ from the original.  This is the
``dice in a cube'' problem: the agent can store what it has learned, but
not the emergent structure of its relationships with other agents.

%% ============================================================
\section{Objections}
\label{sec:objections}

\subsection{``This is just multi-agent systems.''}

Multi-agent systems have existed for decades
(Section~\ref{sec:related}).  The novelty is not the architecture; it is
the nature of the agents.  Traditional multi-agent systems use agents
that follow rules, match patterns, and execute programmed behaviors.
The constraint architecture in such systems must compensate for the
agents' lack of understanding.  In Emergence, contextual understanding
is a first-class capability.  The constraint architecture does not
compensate for a deficit; it channels a surplus.

\subsection{``You cannot rely on LLM judgment.''}

If LLM judgment is unreliable in all contexts, then LLMs are useless.
If it is reliable in some contexts and unreliable in others, the
engineering question is how to create the contexts in which it is
reliable.  That is what the constraint architecture does: schemas ensure
structural validity, scoring functions calibrate judgments, consensus
mechanisms prevent single-agent errors from propagating, and lifecycle
rules ensure that consistently wrong agents decay.  The paradigm does
not assume perfect agents.  It assumes \emph{bounded} agents operating
within a structure that compensates for individual limitations---the
same assumption that underpins every functioning human organization.

\subsection{``Emergent behavior is unpredictable.''}

Yes.  That is the point.  If behavior were fully predictable, the system
would not be discovering anything---it would be computing something the
programmer already knew.  The question is not whether behavior is
predictable but whether it is \emph{bounded}---whether the constraint
architecture prevents pathological behavior while permitting useful
surprise.  This is the same question governing any adaptive system, from
immune responses to market economies to scientific communities.

\subsection{``This will be commoditized into frameworks.''}

It will.  And the commoditization will produce a useful litmus test.
When the paradigm attracts ``emergence frameworks'' that are
orchestration with randomness, the name itself provides the criterion:
\emph{did behavior emerge, or was it scripted?}  If you can read the
source and predict the output, it is not Emergence.  It is procedural
programming with extra steps.  A system in which the programmer
specified the behavior, however indirectly, is not exhibiting emergence.
A system in which the programmer specified the constraints, and behavior
arose from agents operating within those constraints, is.

%% ============================================================
\section{Conclusion}
\label{sec:conclusion}

Programming paradigms arrive when a new capability makes a new
relationship between programmer and machine possible.  Structured
programming arrived when hardware supported abstraction.  Object-oriented
programming arrived when GUIs demanded encapsulation.  Functional
programming arrived (the second time) when multicore demanded
immutability.

Emergence arrives because large language models have created
computational units with contextual understanding.  For the first time,
the receiver of a message can decide what it means.  For the first
time, the computational primitive can exercise judgment.  For the first
time, the programmer can specify \emph{what matters} and trust the
system to figure out \emph{how to assess it.}

This was Kay's vision.  Not classes and inheritance.  Not design
patterns and factories.  Autonomous entities interpreting messages,
coordinating through shared environments, producing behavior that
emerges from their interaction rather than from central specification.

The programmer does not write the program.  The programmer writes the
conditions under which the program emerges.  The constraints are the
artifact.  The computation is what arises within them.

The name carries an honest admission.  Emergence is not deterministic.
The same constraints, the same agents, the same signals will not produce
identical trajectories.  They will produce structurally consistent
trajectories---bounded by the same constraints, shaped by the same
scoring functions, governed by the same lifecycle rules---but the
specific paths will vary.  This is the admission that distinguishes the
paradigm from its predecessors and the commitment that makes it
engineering: design the constraints well enough that every trajectory is
acceptable, and the non-determinism becomes a feature.  Design them
poorly, and it becomes noise.

That is the skill.  That is the discipline.  And for the first time, the
technology exists to practice it.

\bibliographystyle{plainnat}
\bibliography{emergence}

\end{document}
